\documentclass[10pt,a4paper]{amsart}
\usepackage[margin=19mm]{geometry}
\usepackage{amsmath,amssymb,mathtools}
\usepackage[scr=rsfs,cal=euler]{mathalfa}
\usepackage{xfrac}
\usepackage[unicode,hidelinks,bookmarks=false]{hyperref}
\newcommand{\ie}{i.e.\ }
\newcommand{\cf}{cf.\ }
\newcommand{\resp}{respectively\ }
\newcommand{\Subsection}{\subsection}
\providecommand{\nobreakdash}{\nobreak\hbox{-}}
\setlength{\emergencystretch}{2em}
\multlinegap0pt
\usepackage{amssymb}
\usepackage{enumerate}
\usepackage[normalem]{ulem}
\newtheorem{lemma}{Lemma}
\newcommand{\C}{\mathbb{C}}
\title{A Matsushima theorem for K-polystable polarised smooth Fano threefolds}
\author{Hamid Abban, Paolo Cascini, Ivan Cheltsov}
\date{Source: arXiv:2604.20440v1 (CC BY 4.0)}
\begin{document}
\maketitle

\noindent\emph{Source and licence.} A Matsushima theorem for K-polystable polarised smooth Fano threefolds. Version arXiv:2604.20440v1 (CC BY 4.0), distributed by arXiv under the Creative Commons Attribution 4.0 International licence, http://creativecommons.org/licenses/by/4.0/, as recorded in the arXiv licence metadata for this exact version and shown on the abstract page https://arxiv.org/abs/2604.20440v1. Full licence notice: \texttt{licenses/arxiv-2604.20440-CC-BY-4.0.txt}.\par\smallskip

\noindent\textit{Editorial scope.} Counted unit: the article's lemma lemma:3-16-DF (source0.tex lines 668--680). The article's complete original proof of it is delivered with it (source0.tex lines 682--760, one \texttt{proof} environment of 79 lines). No mathematics is written, repaired or re-derived: the statement and the proof are the article's own lines, in the article's own notation, and no retained line is substituted or re-flowed.  Retained uncounted as the source-local context the proof needs: lines 367--374; lines 537--544.  Nothing was omitted from any retained span, so the package carries no omission record.  No retained line contains a citation, so the wrapper carries no bibliography and no bibliography entry.  No retained line contains a figure environment, an image-inclusion command or drawing code, so no image is delivered and the package declares no asset dependency.  Editorial formatting only, in addition: an AMS \texttt{amsart} preamble that restates the article's own declarations and macros used by the retained lines, namely the environments \texttt{lemma} on separate counters, together with the standard packages those lines need.  The retained environments print in this document as Lemma 1 in order of appearance, whereas the article prints the same items with the numbering of its own full document.  The complete native source of the article as distributed in the arXiv e-print for this exact version is bundled unchanged as \texttt{sources/198935/source0.tex}.
\begin{equation}
\label{equation:ab}
b_0=\sum_{P\in X^{\C^*}}\frac{\mu^4_L(P)}{24\,\alpha_{1}(P)\alpha_{2}(P)\alpha_{3}(P)}
\quad
\text{and}
\quad
b_1=\sum_{P\in X^{\C^*}}\frac{\mu^3_L(P)(\alpha_{1}(P)+\alpha_{2}(P)+\alpha_{3}(P))}{12\,\alpha_{1}(P)\alpha_{2}(P)\alpha_{3}(P)}.
\end{equation}

\begin{lemma}
\label{lemma:2-26-DF}
For the product test configuration that is induced by the described $\C^*$-action, the~Donaldson--Futaki invariant of $(X,L)$ is
\begin{equation}
\label{equation:2-26-DF}
-\frac{ab\bigl(6a^4+35a^3b+80a^2b^2+75ab^3+25b^4\bigr)}{4(a+b)(2a^2+10ab+5b^2)}.
\end{equation}
\end{lemma}

\begin{lemma}
\label{lemma:3-16-DF}
For the product test configuration induced by the described $\C^*$-action on $X$, the~Donaldson--Futaki invariant of $(X,L)$ is
$$
-\frac{f(a,b,c)}{4(a^{3} + 3 a^{2} b + 6 a^{2} c + 3 a b^{2} + 12 a b c + 3 a c^{2} + 3 b^{2} c + 3 b c^{2})},
$$
where
\begin{multline*}
f(a,b,c)=3a^{4} b^{2} + 6 a^{4} b c + 6 a^{4} c^{2} + 4 a^{3} b^{3} + 25 a^{3} b^{2} c + 67 a^{3} b c^{2} + 48 a^{3} c^{3}+\\
+ 24 a^{2} b^{3} c + 129 a^{2} b^{2} c^{2} + 204 a^{2} b c^{3} + 90 a^{2} c^{4}+ 12 a b^{4} c + 102 a b^{3} c^{2} + \\
+243 a b^{2} c^{3} + 201 a b c^{4} + 36 a c^{5}+ 12 b^{4} c^{2} + 54 b^{3} c^{3} + 78 b^{2} c^{4} + 36 b c^{5}.
\end{multline*}
\end{lemma}

\begin{proof}
The proof is very similar to  the proof of Lemma~\ref{lemma:2-26-DF}. For convenience of the reader, we present some details.
The Donaldson--Futaki invariant of $(X,L)$ is $\frac{b_0a_1-b_1a_0}{a_0}$, where
$$
a_0=\frac{L^3}{6}=\frac{a^3+3a^2b+6a^2c+3ab^2+12abc+3ac^2+3b^2c+3bc^2}{6},
$$
$$
a_1=\frac{(-K_X\cdot L^2)}{4}=\frac{2(2a^2+4ab+5ac+b^2+4bc+c^2)}{4},
$$
and $b_0$ and $b_1$ can be computed using  \eqref{equation:ab} as follows.

Let $X^{\C^*}$ be the set of $\mathbb{C}^*$-fixed points in $X$.
Then $X^{\C^*}$ consists of eight points.
To describe them, observe that our $\mathbb{C}^*$-action fixes exactly $4$ points in $\mathbb{P}^3$. These points are
$$
O=[1:0:0:0],\,
Q_1=[0:1:0:0],\,
Q_2=[0:0:1:0],\,
Q_3=[0:0:0:1].
$$
The twisted cubic $C_3$ contains $O$ and $Q_3$, and it does not contain $Q_1$ and $Q_2$.
We denote by $\bar{Q}_1$, $\bar{Q}_2$, $\bar{Q}_3$ the preimages in $V_7$ of the points $Q_1$, $Q_2$, $Q_3$, respectively.
Set $\overline{E}=\pi(E)$. Then $\overline{E}\simeq\mathbb{P}^2$ is $\phi$-exceptional divisor,
and our $\C^*$-action acts on $\overline{E}$ with weights $(1,2,3)$, so it fixes three points,
which we denote by $\bar{Q}_4$, $\bar{Q}_5$, $\bar{Q}_6$.
Note that one of these points is contained in $C$. Without loss of generality, we may assume that $\bar{Q}_4\in C$.
Thus, $C$ contains $\bar{Q}_3$ and $\bar{Q}_4$, and it does not contain  $\bar{Q}_1$, $\bar{Q}_2$, $\bar{Q}_5$, $\bar{Q}_6$.

We denote by $P_1$, $P_2$, $P_5$ and $P_6$ the preimages on $X$ of the points $\bar{Q}_1$, $\bar{Q}_2$, $\bar{Q}_5$, $\bar{Q}_6$, respectively.
Let $\ell_3$ and $\ell_4$ be the fibres of the natural projection $F\to C$ over the points $\bar{Q}_3$ and $\bar{Q}_4$, respectively.
Then $\ell_3$ contains two $\mathbb{C}^*$-fixed points, which we denote by $P_3$ and $P_3^\prime$.
Similarly, the curve $\ell_4$ contains two $\mathbb{C}^*$-fixed points, which we denote by $P_4$ and $P_4^\prime$.
Then
$$
X^{\C^*}=\big\{P_1,P_2,P_3,P_3^\prime,P_4,P_4^\prime,P_5,P_6\big\}.
$$
By construction, $P_1$ and $P_2$ are not contained in $E\cup F$, while  $P_3$ and $P_3^\prime$ are contained in $F$,
but $P_3$ and $P_3^\prime$ are not contained in $E$.
Similarly, points $P_4$ and $P_4^\prime$ are contained in $\ell_4=E\cap F$, while points $P_5$ and $P_6$ are contained in $E$,
but they are not contained in $F$.

For every $P\in X^{\C^*}$, the group $\mathbb{C}^*$ faithfully acts on the cotangent space $T_P^*(X)$,
and we denote by $\alpha_1(P)$, $\alpha_2(P)$, $\alpha_3(P)$ the weights of this action.
Possibly after swapping $P_3\leftrightarrow P_3^\prime$, $P_4\leftrightarrow P_4^\prime$, $P_5\leftrightarrow P_6$, 
the weights of these eight $\mathbb{C}^*$-actions are computed as follows:
\begin{align*}
\big(\alpha_1(P_1),\alpha_2(P_1),\alpha_3(P_1)\big)&=(-1,1,2),&\quad
\big(\alpha_1(P_2),\alpha_2(P_2),\alpha_3(P_2)\big)&=(-2,-1,1),&\quad\\
\big(\alpha_1(P_3),\alpha_2(P_3),\alpha_3(P_3)\big)&=(-2,-1,-1),&\quad
\big(\alpha_1(P_3^\prime),\alpha_2(P_3^\prime),\alpha_3(P_3^\prime)\big)&=(1,-3,-1),\\
\big(\alpha_1(P_4),\alpha_2(P_4),\alpha_3(P_4)\big)&=(-1,2,1),&\quad
\big(\alpha_1(P_4^\prime),\alpha_2(P_4^\prime),\alpha_3(P_4^\prime)\big)&=(1,1,1),\\
\big(\alpha_1(P_5),\alpha_2(P_5),\alpha_3(P_5)\big)&=(-1,2,1),&\quad
\big(\alpha_1(P_6),\alpha_2(P_6),\alpha_3(P_6)\big)&=(-2,-1,3),&\quad
\end{align*}
Similarly, we compute the fibre weights for  $H$, $E$ and $F$ as follows:
\begin{multline*}
\quad \quad \quad \quad \quad \mu_H(P_1)=1, \mu_H(P_2)=2, \mu_H(P_3)=3, \mu_H(P_3^\prime)=3, \\
\mu_H(P_4)=0, \mu_H(P_4^\prime)=0, \mu_H(P_5)=0, \mu_H(P_6)=0,\quad \\
\mu_E(P_1)=0, \mu_E(P_2)=0, \mu_E(P_3)=0, \mu_E(P_3^\prime)=0, \\
\mu_E(P_4)=-1, \mu_E(P_4^\prime)=-1, \mu_E(P_5)=-2, \mu_E(P_6)=-3, \\
\mu_F(P_1)=0, \mu_F(P_2)=0, \mu_F(P_3)=2, \mu_F(P_3^\prime)=3, \\
\mu_F(P_4)=-2, \mu_F(P_4^\prime)=-1,  \mu_F(P_5)=0, \mu_F(P_6)=0.\quad \quad \quad \quad \quad
\end{multline*}
Thus, the fibre weight for the line bundle $L$ is
\begin{multline*}
\mu_L(P_1)=a+b+2c,\quad
\mu_L(P_2)=2a+2b+4c,\quad
\mu_L(P_3)=3a+3b+4c,\quad
\mu_L(P_3^\prime)=3a+3b+3c,\quad\\
\mu_L(P_4)=b+3c,\quad
\mu_L(P_4^\prime)=b+2c,\quad
\mu_L(P_5)=2b+2c,\quad
\mu_L(P_6)=3b+3c,\quad
\end{multline*}

Now, using \eqref{equation:ab}, we compute $b_0$ and $b_1$ in $\frac{b_0a_1-b_1a_0}{a_0}$, which results
in the required formula for the Donaldson--Futaki invariant.
\end{proof}

\end{document}
